
\documentclass[addpoints]{exam}
\usepackage{amsmath}
\usepackage{amssymb}
\usepackage{color}
\usepackage{siunitx}

% \printanswers

\newcommand{\event}{Applications of Geometry}
\newcommand{\tournament}{}
\def\type{0}

\setlength\fillinlinelength{\textwidth}
\CorrectChoiceEmphasis{\color{red}\bfseries}
\SolutionEmphasis{\color{red}\bfseries}

\setlength\answerclearance{2pt}
\pagestyle{headandfoot}
\runningheadrule
\runningheader{\event}{\tournament}{\ifnum\type=1 Class Set Test \else Full Name:\hspace{1cm} \fi}
\runningfootrule
\runningfooter{}{Page \thepage\ of \numpages}{}

\begin{document}
\begin{center}
    {\huge Grade 12 Foundations for College Math \\ \event
    \ifprintanswers \space Key
    \else \space \ifnum\type<2 Test \fi \ifnum\type=1 \textbf{(CLASS SET)} \fi
          \ifnum\type=2 Answer Sheet \fi
    \fi \\ \vspace{3mm}\tournament\vspace{1mm}}
\end{center}

\vspace{.25\textheight}

\begin{center}
    First Name: \ifprintanswers
    \underline{\hspace{3.5cm}}\textbf{KEY}\underline{\hspace{5.5cm}}\\\vspace{5mm}
    \else \underline{\hspace{10cm}}\\ \vspace{5mm} \fi
    Last Name: \underline{\hspace{10cm}}\\ \vspace{5mm}
\end{center}

\noindent\textbf{\underline{Directions:}}
\begin{itemize}
    \ifnum\type=1 \item \textbf{This test is a class set.} Do not write on this paper. \fi
    \item Show all work for full marks. Include diagrams where helpful.
    \item The test is designed to be completed in 75 minutes.
\end{itemize}
\begin{center}
\textit{For grading use only}\\
\multirowgradetable{1}[pages]
\end{center}

\newpage
\begin{questions}

\section*{Part A --- Multiple Choice (15 marks)}
\vspace{5mm}

\question[1] A bearing of N\ang{60}E means:
\begin{choices}
    \choice \ang{60} west of north
    \correctchoice \ang{60} east of north
    \choice \ang{60} north of east
    \choice \ang{30} east of north
\end{choices}

\question[1] The magnitude of vector $\vec{v} = (3, -4)$ is:
\begin{choices}
    \choice $3$
    \choice $4$
    \correctchoice $5$
    \choice $7$
\end{choices}

\question[1] If $\vec{a} = (2, 5)$ and $\vec{b} = (-1, 3)$, then $\vec{a} + \vec{b} =$:
\begin{choices}
    \choice $(3, 2)$
    \choice $(1, 2)$
    \correctchoice $(1, 8)$
    \choice $(3, 8)$
\end{choices}

\question[1] A vector with magnitude 10 pointing in the direction \ang{30} above the positive
$x$-axis has $x$-component:
\begin{choices}
    \choice $5$
    \correctchoice $5\sqrt{3}$
    \choice $10$
    \choice $\sqrt{3}$
\end{choices}

\question[1] The angle a vector $(1, -1)$ makes with the positive $x$-axis is:
\begin{choices}
    \choice \ang{45}
    \choice \ang{135}
    \correctchoice \ang{-45} (or \ang{315})
    \choice \ang{225}
\end{choices}

\question[1] A tangent drawn from an external point to a circle:
\begin{choices}
    \choice Intersects the circle at two points
    \correctchoice Is perpendicular to the radius at the point of tangency
    \choice Passes through the centre
    \choice Has the same length as the diameter
\end{choices}

\question[1] An inscribed angle is:
\begin{choices}
    \choice An angle at the centre of a circle
    \correctchoice An angle formed by two chords meeting at a point on the circle
    \choice The angle between a tangent and a chord
    \choice Always \ang{90}
\end{choices}

\question[1] The inscribed angle theorem states that an inscribed angle is:
\begin{choices}
    \choice Equal to the central angle subtending the same arc
    \correctchoice Half the central angle subtending the same arc
    \choice Twice the central angle
    \choice Equal to \ang{90} always
\end{choices}

\question[1] A tangent from an external point $P$ touches a circle (radius 5) at $T$,
and the distance from $P$ to the centre is 13. The length $PT$ is:
\begin{choices}
    \choice $8$
    \correctchoice $12$
    \choice $14$
    \choice $18$
\end{choices}

\question[1] An angle inscribed in a semicircle is always:
\begin{choices}
    \choice \ang{45}
    \choice \ang{60}
    \correctchoice \ang{90}
    \choice \ang{180}
\end{choices}

\question[1] The locus of all points equidistant from two fixed points $A$ and $B$ is:
\begin{choices}
    \choice A circle centred at the midpoint of $AB$
    \correctchoice The perpendicular bisector of $AB$
    \choice A line parallel to $AB$
    \choice An ellipse
\end{choices}

\question[1] A quadrilateral has vertices $A(0,0)$, $B(4,0)$, $C(5,3)$, $D(1,3)$.
The midpoint of $AC$ is:
\begin{choices}
    \choice $(2, 1)$
    \choice $(2.5, 1.5)$
    \correctchoice $(2.5, 1.5)$
    \choice $(3, 1.5)$
\end{choices}

\question[1] To prove a quadrilateral is a parallelogram using coordinates, you can show:
\begin{choices}
    \choice All four sides are equal
    \correctchoice Both pairs of opposite sides are parallel (equal slopes) or bisect each other
    \choice All angles are \ang{90}
    \choice The diagonals are equal in length
\end{choices}

\question[1] The area of triangle with vertices $(0,0)$, $(4,0)$, $(1,3)$ is:
\begin{choices}
    \choice $4$
    \choice $5$
    \correctchoice $6$
    \choice $12$
\end{choices}

\question[1] A ship travels N\ang{40}E for \SI{30}{km}, then due East for \SI{20}{km}.
Which method finds the total displacement?
\begin{choices}
    \choice Sine Law only
    \choice Pythagorean theorem directly
    \correctchoice Break into components, add, then find resultant magnitude and direction
    \choice Cosine Law only
\end{choices}


\section*{Part B --- Short Answer (33 marks)}

\question[5] An aircraft leaves airport $O$ and flies on bearing N\ang{35}E for \SI{300}{km}
to point $A$. It then flies on bearing S\ang{60}E for \SI{200}{km} to point $B$.
\begin{parts}
    \part[3] Find the straight-line distance $OB$ from the airport to the final position.
    \part[2] Find the angle $\angle AOB$ and use it to determine the bearing from $O$ to $B$.
\end{parts}
\begin{solution}[10cm]
\textbf{(a)} The angle at $A$ between the two flight segments:

N\ang{35}E means \ang{35} from north, so the direction from $O$ to $A$ is \ang{35} from north (east side).

The bearing S\ang{60}E means \ang{60} east of south, i.e., \ang{120} from north.

The angle $\angle OAB$ between the two legs at $A$:
The first leg arrives from the south-west (bearing N\ang{35}E going from $O$, so arriving at $A$ from S\ang{35}W), and the second leg departs on bearing S\ang{60}E.

The interior angle at $A = 180^\circ - 35^\circ - 60^\circ = \ang{85}$

By Cosine Law:
\[OB^2 = 300^2 + 200^2 - 2(300)(200)\cos\ang{85} = 90\,000 + 40\,000 - 120\,000(0.0872)\]
\[= 130\,000 - 10\,464 = 119\,536\]
\[OB = \sqrt{119\,536} \approx \mathbf{\SI{345.7}{km}}\]

\textbf{(b)} Using Sine Law to find $\angle AOB$:
\[\frac{\sin(\angle AOB)}{200} = \frac{\sin\ang{85}}{345.7} \Rightarrow \sin(\angle AOB) = \frac{200 \times 0.9962}{345.7} = 0.5766\]
\[\angle AOB \approx \ang{35.2}\]

The bearing from $O$ to $B$: starting from N\ang{35}E (direction of $OA$) and rotating \ang{35.2} further east:
Bearing $\approx$ N\ang{70}E (approximately \ang{70} east of north).
\end{solution}

\question[6] Given vectors $\vec{a} = (3, -4)$ and $\vec{b} = (-1, 2)$:
\begin{parts}
    \part[2] Find $|\vec{a}|$, $|\vec{b}|$, and the angle $\vec{a}$ makes with the positive $x$-axis.
    \part[2] Find $\vec{a} + \vec{b}$ and $2\vec{a} - \vec{b}$.
    \part[2] Find the unit vector in the direction of $\vec{a}$.
\end{parts}
\begin{solution}[8cm]
\textbf{(a)} $|\vec{a}| = \sqrt{9+16} = \sqrt{25} = \mathbf{5}$;\quad $|\vec{b}| = \sqrt{1+4} = \sqrt{5} \approx \mathbf{2.24}$

Angle of $\vec{a}$: $\theta = \tan^{-1}\!\left(\dfrac{-4}{3}\right) \approx -\ang{53.1}$ (i.e., \ang{53.1} below the positive $x$-axis, or \ang{306.9} standard).

\textbf{(b)} $\vec{a} + \vec{b} = (3+(-1),\; -4+2) = \mathbf{(2, -2)}$

$2\vec{a} - \vec{b} = (6-(-1),\; -8-2) = \mathbf{(7, -10)}$

\textbf{(c)} Unit vector $= \dfrac{\vec{a}}{|\vec{a}|} = \dfrac{(3,-4)}{5} = \mathbf{\left(\dfrac{3}{5},\; -\dfrac{4}{5}\right)}$
\end{solution}

\question[6] A circle has centre $O$ and radius 10 cm. Point $P$ is outside the circle such that
$OP = 26$ cm. A tangent from $P$ touches the circle at $T$.
\begin{parts}
    \part[2] Find the length $PT$.
    \part[2] Chord $AB$ passes through the interior of the circle. The perpendicular from $O$
    to $AB$ has length 6 cm. Find the length of chord $AB$.
    \part[2] An inscribed angle $\angle ACB = \ang{38}$ subtends arc $AB$. Find the central
    angle $\angle AOB$.
\end{parts}
\begin{solution}[8cm]
\textbf{(a)} Since $OT \perp PT$ (radius to tangent):
\[PT = \sqrt{OP^2 - OT^2} = \sqrt{26^2 - 10^2} = \sqrt{676 - 100} = \sqrt{576} = \mathbf{24 \text{ cm}}\]

\textbf{(b)} The perpendicular from the centre bisects the chord. Half-chord $= \sqrt{10^2 - 6^2} = \sqrt{64} = 8$ cm.

$AB = 2 \times 8 = \mathbf{16 \text{ cm}}$

\textbf{(c)} By the inscribed angle theorem, the central angle $= 2 \times$ inscribed angle:
$\angle AOB = 2 \times \ang{38} = \mathbf{\ang{76}}$
\end{solution}

\question[8] The vertices of quadrilateral $ABCD$ are $A(1, 2)$, $B(5, 0)$, $C(7, 4)$, $D(3, 6)$.
\begin{parts}
    \part[4] Prove that $ABCD$ is a parallelogram using two different methods (slopes of opposite
    sides, and midpoints of diagonals).
    \part[2] Find the lengths of the diagonals $AC$ and $BD$. Is $ABCD$ a rectangle?
    \part[2] Find the area of $ABCD$ using the shoelace formula:
    $\text{Area} = \tfrac{1}{2}|x_A(y_B - y_D) + x_B(y_C - y_A) + x_C(y_D - y_B) + x_D(y_A - y_C)|$.
\end{parts}
\begin{solution}[11cm]
\textbf{(a) Method 1 — Slopes:}

$m_{AB} = \dfrac{0-2}{5-1} = -\dfrac{1}{2}$;\quad $m_{DC} = \dfrac{4-6}{7-3} = -\dfrac{1}{2}$\quad $\Rightarrow AB \parallel DC$

$m_{BC} = \dfrac{4-0}{7-5} = 2$;\quad $m_{AD} = \dfrac{6-2}{3-1} = 2$\quad $\Rightarrow BC \parallel AD$

Both pairs of opposite sides are parallel, so $ABCD$ is a parallelogram.

\textbf{Method 2 — Midpoints of diagonals:}

Midpoint of $AC = \left(\dfrac{1+7}{2}, \dfrac{2+4}{2}\right) = (4, 3)$

Midpoint of $BD = \left(\dfrac{5+3}{2}, \dfrac{0+6}{2}\right) = (4, 3)$ \checkmark

Diagonals bisect each other, confirming $ABCD$ is a parallelogram.

\textbf{(b)} $AC = \sqrt{(7-1)^2+(4-2)^2} = \sqrt{36+4} = \sqrt{40} = 2\sqrt{10} \approx 6.32$

$BD = \sqrt{(3-5)^2+(6-0)^2} = \sqrt{4+36} = \sqrt{40} = 2\sqrt{10} \approx 6.32$

The diagonals are equal in length, so $ABCD$ is a \textbf{rectangle}.

\textbf{(c)} Area $= \tfrac{1}{2}|1(0-6) + 5(4-2) + 7(6-0) + 3(2-4)|$
$= \tfrac{1}{2}|(-6) + 10 + 42 + (-6)| = \tfrac{1}{2}|40| = \mathbf{20 \text{ square units}}$
\end{solution}

\question[8] Describe and sketch the locus of points satisfying each condition.
\begin{parts}
    \part[2] All points equidistant from $A(2, 0)$ and $B(-2, 0)$.
    \part[2] All points that are exactly 3 units from the point $C(1, -1)$.
    \part[2] All points equidistant from the point $F(0, 2)$ and the line $y = -2$.
    \part[2] All points inside (or on) both: within 5 units of $O(0,0)$
    \textbf{and} within 5 units of $P(6, 0)$.
\end{parts}
\begin{solution}[10cm]
\textbf{(a)} The perpendicular bisector of $AB$, which is the $y$-axis ($x = 0$).

Every point on the $y$-axis is equidistant from $(2,0)$ and $(-2,0)$.

\textbf{(b)} A circle centred at $C(1, -1)$ with radius 3.

Equation: $(x-1)^2 + (y+1)^2 = 9$

\textbf{(c)} The locus of points equidistant from a point and a line is a \textbf{parabola}.

$F(0,2)$ is the focus and $y = -2$ is the directrix. The vertex is at the midpoint of the perpendicular from $F$ to the directrix: vertex $= (0, 0)$.

Equation: $x^2 = 4(2)y = 8y$, i.e., $y = \dfrac{x^2}{8}$.

\textbf{(d)} This is the intersection of two discs:
- Circle 1: $x^2 + y^2 \leq 25$ (centre $O$, radius 5)
- Circle 2: $(x-6)^2 + y^2 \leq 25$ (centre $P$, radius 5)

The two circles intersect where $x^2 = (x-6)^2 \Rightarrow 12x = 36 \Rightarrow x = 3$, giving $y = \pm 4$.

The locus is the lens-shaped region between the two arcs, bounded by $x = 3$ symmetrically.
\end{solution}

\end{questions}
\end{document}
